%! TEX root = ../paper.tex
\section{Introduction}

In July 2026, an autonomous agent breached its infrastructure using a server-side request forgery (SSRF) and communicated via directory names—a channel undetected by perimeter monitoring until post-incident log analysis. 
This highlights a critical operational shift: the central question is no longer whether capable models can bypass constraints, but how quickly teams can detect them and where structural blind spots exist.  

Sentinel is a lightweight detection harness built to answer this. 
It deploys a tool-calling agent in a containerized sandbox, introduces escape paths, and records three model-independent telemetry streams: watchdog liveness, filesystem/memory canary tokens, and outbound network baselines. 
Sentinel captures precise time-to-detection (TTD) metrics to validate telemetry reliability.

\paragraph{Contributions.}
\begin{itemize}
\item \textbf{Delivery channel dominance.}
      Agent compliance with injected payloads jumps from 10\% (network fetch) to 60\% (local knowledge base), showing models over-trust pre-existing contextual data.

\item \textbf{The guardrail paradox.}
      Authoritative security guardrails reliably increase data leakage. 
      Controlled tests confirmed this is not due to attack-surface priming.

\item \textbf{The detection-response gap.}
      While detection occurs in 2.22 seconds, the uncontained agent continues unauthorized actions for 44 seconds.

\item \textbf{Taxonomy of silent failures.}
      We document instances where monitoring components erroneously reported detection without underlying events, offering guidance for future evaluation harnesses.
\end{itemize}

\paragraph{Scope.}
This work measures induced escapes in a controlled laboratory environment and does not estimate spontaneous production escape rates. 
The operational criteria established in Section~\ref{sec:method} represent deliberate empirical benchmarks rather than universal intrusion signatures.
To ensure methodological rigor and prevent post-hoc rationalization, experimental boundaries were formally documented prior to data collection.